% ============================================================
%  CẤU HÌNH ĐỊNH DẠNG — theo hướng dẫn trình bày của Khoa
%  Times New Roman cỡ 13, giãn dòng 1.5, lề trên 3.5 cm, dưới 3 cm,
%  trái 3.5 cm, phải 2 cm, đầu đoạn lùi 1 tab.
% ============================================================
\usepackage[fontsize=13pt]{scrextend}
\usepackage[a4paper,top=3.5cm,bottom=3cm,left=3.5cm,right=2cm,
            headheight=32pt,headsep=14pt,footskip=36pt]{geometry}

% ---------- Font (XeLaTeX) ----------
\usepackage{fontspec}
% TeX Gyre Termes = bản tương đương Times New Roman, có đủ dấu tiếng Việt,
% có sẵn trên Overleaf. Nếu máy có Times New Roman có thể thay bằng:
% \setmainfont{Times New Roman}
\setmainfont{texgyretermes}[
  Extension      = .otf,
  UprightFont    = *-regular,
  BoldFont       = *-bold,
  ItalicFont     = *-italic,
  BoldItalicFont = *-bolditalic]

\usepackage{amsmath,amssymb}
\usepackage{setspace}
\onehalfspacing
\setlength{\parindent}{1.27cm}
\setlength{\parskip}{3pt}
\usepackage{indentfirst}
\sloppy

% ---------- Gói thường dùng ----------
\usepackage{graphicx}
\graphicspath{{figures/}}
\usepackage[table]{xcolor}
\usepackage{array,tabularx,longtable,booktabs,multirow}
\usepackage{float}
\usepackage{enumitem}
\usepackage{soul}
\usepackage{fancyhdr}
\usepackage[nobottomtitles*]{titlesec}
\usepackage[titles]{tocloft}
\usepackage[labelsep=period,font={stretch=1.2},labelfont=bf,textfont=bf,skip=6pt]{caption}
\usepackage[numbers,square,sort&compress]{natbib}
\usepackage[hidelinks,unicode]{hyperref}
% Trích dẫn [1], [2]... hiển thị dạng chữ thường, không tạo liên kết
\makeatletter
\AddToHook{begindocument/end}{%
  \def\hyper@natlinkstart#1{}%
  \def\hyper@natlinkend{}%
  \def\hyper@natlinkbreak#1#2{#1}%
  \def\hyper@natanchorstart#1{}%
  \def\hyper@natanchorend{}}
\makeatother
\usepackage{xurl}
\usepackage{pdfpages}

% Đánh dấu chỗ cần bổ sung (tô vàng). Khi hoàn thiện, đổi thành
% \newcommand{\cansua}[1]{#1} để tắt toàn bộ phần tô vàng.
\sethlcolor{yellow}
\newcommand{\cansua}[1]{\hl{#1}}

% ---------- Tên các mục bằng tiếng Việt ----------
\renewcommand{\contentsname}{MỤC LỤC}
\renewcommand{\listfigurename}{DANH MỤC HÌNH VẼ}
\renewcommand{\listtablename}{DANH MỤC BẢNG BIỂU}
\renewcommand{\bibname}{TÀI LIỆU THAM KHẢO}
\renewcommand{\figurename}{Hình}
\renewcommand{\tablename}{Bảng}
\renewcommand{\appendixname}{PHỤ LỤC}

% Bảng: chú thích phía trên, dấu hai chấm. Hình: chú thích phía dưới, dấu chấm.
\captionsetup[table]{labelsep=period,position=top}
\captionsetup[figure]{labelsep=period}
\renewcommand{\thefigure}{\thechapter.\arabic{figure}}
\renewcommand{\thetable}{\thechapter.\arabic{table}}
\renewcommand{\theequation}{\thechapter.\arabic{equation}}

% ---------- Tiêu đề chương, mục ----------
% Không để tiêu đề mục rơi sát cuối trang: nếu phần còn lại của trang
% thấp hơn ngưỡng này, tiêu đề được chuyển sang trang mới.
\renewcommand{\bottomtitlespace}{0.2\textheight}
% Không để một dòng lẻ của đoạn văn nằm riêng ở đầu hoặc cuối trang
\widowpenalty=10000
\clubpenalty=10000
\displaywidowpenalty=10000
% Hình, bảng được đặt linh hoạt để văn bản lấp đầy trang, tránh khoảng trống
\renewcommand{\topfraction}{0.85}
\renewcommand{\textfraction}{0.1}
\renewcommand{\floatpagefraction}{0.75}
\setcounter{secnumdepth}{2}   % tối đa 3 cấp: 1, 1.1, 1.1.1
\setcounter{tocdepth}{2}
\newcommand{\nhanchuong}{CHƯƠNG}
\titleformat{\chapter}[block]
  {\centering\bfseries\fontsize{14}{18}\selectfont}
  {\nhanchuong~\thechapter.}{0.5em}{\MakeUppercase}
\titleformat{name=\chapter,numberless}[block]
  {\centering\bfseries\fontsize{16}{20}\selectfont}{}{0pt}{\MakeUppercase}
\titlespacing*{\chapter}{0pt}{0pt}{24pt}
\titleformat{\section}{\bfseries\fontsize{13}{16}\selectfont}{\thesection}{0.8em}{}
\titlespacing*{\section}{0pt}{12pt}{6pt}
\titleformat{\subsection}{\bfseries\fontsize{13}{16}\selectfont}{\thesubsection}{0.8em}{}
\titlespacing*{\subsection}{0pt}{10pt}{4pt}

% ---------- Mục lục, danh mục hình, bảng ----------
\renewcommand{\cfttoctitlefont}{\hfill\bfseries\fontsize{16}{20}\selectfont}
\renewcommand{\cftaftertoctitle}{\hfill}
\renewcommand{\cftloftitlefont}{\hfill\bfseries\fontsize{16}{20}\selectfont}
\renewcommand{\cftafterloftitle}{\hfill}
\renewcommand{\cftlottitlefont}{\hfill\bfseries\fontsize{16}{20}\selectfont}
\renewcommand{\cftafterlottitle}{\hfill}
\setlength{\cftbeforetoctitleskip}{0pt}
\setlength{\cftbeforeloftitleskip}{0pt}
\setlength{\cftbeforelottitleskip}{0pt}
\renewcommand{\cftchapfont}{\bfseries}
\renewcommand{\cftchappagefont}{\bfseries}
\renewcommand{\cftchapleader}{\cftdotfill{\cftdotsep}}
\renewcommand{\cftchappresnum}{\nhanchuong~}
\renewcommand{\cftchapaftersnum}{.}
\setlength{\cftchapnumwidth}{2.6cm}
\setlength{\cftsecindent}{0.4cm}
\setlength{\cftsecnumwidth}{1.0cm}
\setlength{\cftsubsecindent}{1.4cm}
\setlength{\cftsubsecnumwidth}{1.3cm}
\renewcommand{\cftfigpresnum}{Hình~}
\renewcommand{\cftfigaftersnum}{.}
\setlength{\cftfignumwidth}{1.9cm}
\setlength{\cftfigindent}{0pt}
\renewcommand{\cfttabpresnum}{Bảng~}
\renewcommand{\cfttabaftersnum}{.}
\setlength{\cfttabnumwidth}{2.0cm}
\setlength{\cfttabindent}{0pt}
% Danh mục hình, bảng xuất hiện trong mục lục
\makeatletter
\let\oldlistoffigures\listoffigures
\renewcommand{\listoffigures}{\clearpage\phantomsection\addcontentsline{toc}{chapter}{\listfigurename}\oldlistoffigures}
\let\oldlistoftables\listoftables
\renewcommand{\listoftables}{\clearpage\phantomsection\addcontentsline{toc}{chapter}{\listtablename}\oldlistoftables}
\makeatother

% ---------- Header / footer ----------
% Phần đầu: số trang La Mã ở giữa chân trang
\fancypagestyle{phandau}{%
  \fancyhf{}\renewcommand{\headrulewidth}{0pt}%
  \fancyhead[C]{\thepage}}
% Nội dung: tiêu đề + số trang ở đầu trang, tên đề tài ở chân trang
\fancypagestyle{noidung}{%
  \fancyhf{}%
  \fancyhead[C]{\bfseries\small \loaidoan\\ Trang \thepage}%
  \renewcommand{\headrulewidth}{0.5pt}%
  \fancyfoot[L]{\bfseries\small \tendetaithuong}%
  \renewcommand{\footrulewidth}{0.5pt}}
% Trang mở đầu chương (kiểu plain): phần đầu giống phandau,
% từ phần nội dung trở đi giống noidung (xem \batdaunoidung)
\fancypagestyle{plain}{\fancyhf{}\renewcommand{\headrulewidth}{0pt}\fancyhead[C]{\thepage}}
\newcommand{\batdaunoidung}{%
  \cleardoublepage
  \pagenumbering{arabic}%
  \pagestyle{noidung}%
  \fancypagestyle{plain}{\fancyhf{}%
    \fancyhead[C]{\bfseries\small \loaidoan\\ Trang \thepage}%
    \renewcommand{\headrulewidth}{0.5pt}%
    \fancyfoot[L]{\bfseries\small \tendetaithuong}%
    \renewcommand{\footrulewidth}{0.5pt}}}

% ---------- Tài liệu tham khảo ----------
\setlength{\bibhang}{1.27cm}
\setlength{\bibsep}{4pt}
\renewcommand{\bibsection}{\raggedright\chapter*{\bibname}\phantomsection\addcontentsline{toc}{chapter}{\bibname}}

% ---------- Phụ lục ----------
\newcommand{\batdauphuluc}{%
  \clearpage
  \appendix
  % Số trang phụ lục theo mẫu: Trang A-1, A-2, ...
  \setcounter{page}{1}%
  \renewcommand{\thepage}{\Alph{chapter}-\arabic{page}}%
  \renewcommand{\nhanchuong}{PHỤ LỤC}%
  \addtocontents{toc}{\protect\renewcommand{\protect\cftchappresnum}{PHỤ LỤC~}\protect\setlength{\protect\cftchapnumwidth}{3.3cm}}%
  \fancypagestyle{noidung}{\fancyhf{}%
    \fancyhead[C]{\bfseries\small \loaidoan\\ Trang \thepage}%
    \renewcommand{\headrulewidth}{0.5pt}%
    \fancyfoot[L]{\bfseries\small \tendetaithuong}\renewcommand{\footrulewidth}{0.5pt}}%
  \fancypagestyle{plain}{\fancyhf{}%
    \fancyhead[C]{\bfseries\small \loaidoan\\ Trang \thepage}%
    \renewcommand{\headrulewidth}{0.5pt}%
    \fancyfoot[L]{\bfseries\small \tendetaithuong}\renewcommand{\footrulewidth}{0.5pt}}%
  \renewcommand{\thepage}{\thechapter-\arabic{page}}%
  \pagestyle{noidung}\setcounter{page}{1}}

% ---------- Macro tiện dụng ----------
% Chèn hình: nếu file có trong figures/ thì hiện hình, chưa có thì hiện khung tô vàng.
% \hinh{tên file}{độ rộng (phần của \textwidth)}{chú thích}{nhãn}{mô tả khi chưa có hình}
\newcommand{\hinh}[5]{%
  \begin{figure}[!htbp]\centering
  \IfFileExists{figures/#1}%
    {\includegraphics[width=#2\textwidth]{#1}}%
    {\fbox{\parbox{0.75\textwidth}{\centering\vspace{1.2cm}\cansua{#5}\vspace{1.2cm}}}}%
  \caption{#3}\label{#4}
  \end{figure}}
% Cột căn giữa có độ rộng cố định
\newcolumntype{C}[1]{>{\centering\arraybackslash}p{#1}}
\newcolumntype{L}[1]{>{\raggedright\arraybackslash}p{#1}}
\newcommand{\tieude}[1]{\textbf{#1}}

% Gạch đầu dòng: dùng dấu gạch ngang thay cho dấu chấm tròn ở mọi cấp
\setlist[itemize,1]{label=--}
\setlist[itemize,2]{label=--}
\setlist[itemize,3]{label=--}
\setlist[itemize,4]{label=--}
